% Worked operation using existing chapter styles.
\begin{utpcompactlabfigure}[At a matching node: return the subtree that replaces it]
\begin{labgraphic}
\begin{tikzpicture}[font=\small]

\node[utpstep,align=center] (q) at (5,0) {\texttt{value == root->data}};
\node[utpstep,align=left,text width=4.8cm] (a) at (1,-2) {\textbf{Zero or one child}\\Choose the existing child (or null).\\Delete root; return child.};
\node[utpstep,align=left,text width=5.3cm] (b) at (9,-2) {\textbf{Two children}\\Find leftmost node of right subtree.\\Copy its data; delete it below.\\Return the original root.};
\draw[utparrow] (q)--node[utplabel,above left,align=center]{left is null\\or right is null}(a);
\draw[utparrow] (q)--node[utplabel,above right]{both children exist}(b);
\node[utplabel,align=center] at (5,-3.5) {The caller assigns the returned pointer to its child link (or to the tree root).};
\end{tikzpicture}
\end{labgraphic}
\end{utpcompactlabfigure}
